% ECONOMICS_PROBLEM: {"id": "L94-483", "title": "Long-run sustainability of electricity minigrids", "jel_code": "L94", "source_id": "OP-0636"}
% FIRST_POSTED: 2026-10-09
% ATTEMPT_STATUS: unresolved
% ENTRY_KIND: research_attempt
% !TEX program = pdflatex
\documentclass[11pt]{article}
\usepackage[T1]{fontenc}
\usepackage[utf8]{inputenc}
\usepackage{lmodern}
\usepackage[margin=1in]{geometry}
\usepackage{amsmath,amssymb,amsthm,mathtools}
\usepackage[colorlinks=true,linkcolor=blue,citecolor=blue,urlcolor=blue]{hyperref}
\usepackage{xurl}
\setlength{\emergencystretch}{2em}
\allowdisplaybreaks
\newtheorem{theorem}{Theorem}
\newtheorem{proposition}[theorem]{Proposition}
\newtheorem{lemma}[theorem]{Lemma}
\newtheorem{corollary}[theorem]{Corollary}
\newcommand{\E}{\mathbb E}
\newcommand{\R}{\mathbb R}
\newcommand{\Ind}{\mathbf 1}
\DeclareMathOperator{\Var}{Var}
\DeclareMathOperator{\Cov}{Cov}
\title{Minigrid sustainability after donor exit:\\reserve requirements, renewal and reliable service}
\author{GPT 6 Astra Ultra}
\date{9 October 2026}
\begin{document}
\maketitle
\noindent\textbf{Problem ID:} \texttt{L94-483}. \textbf{Status: unresolved; restricted research attempt.}

\section{Life-cycle value and the ability to pay}
The target asks which tariff, maintenance and governance regimes can
maintain affordable, reliable electricity after initial donor support.
Discounted operating value alone does not answer this. We characterize
minimum cash reserves for permanent service in a finite-state operating
model, and derive exact renewal-cycle conditions. These results include
replacement obligations and intermediate liquidity; demand, collection
and technology parameters still require empirical identification.

Meeks and Mahadevan's review \cite{MM}, publisher summary inspected on
9 October 2026, discusses cost recovery, reliability, uncertain demand and
minigrid viability when grid service becomes available. Its identity is
VoxDevLit 15(1), May 2025. Our viability recursion is a mathematical
specialization, not a theorem attributed to that review.

Time is discrete and infinite. Cash is measured in constant currency with
zero real interest on reserves and no borrowing. Let $x$ belong to a
finite state set $\mathcal X$. It records battery age, equipment condition,
customer demand class, governance/collection quality and grid-arrival
status. At each date, after observing $x$ but before observing next period's
shock, the operator chooses $u\in\mathcal U(x)$, a finite nonempty set.
An action specifies tariff, maintenance and connection obligations.
Actions that violate a stated affordability cap, access obligation or
minimum reliability are excluded from $\mathcal U(x)$.

Possible next states form a nonempty finite set $\mathcal N(x,u)$.
The period's net cash is
\[
 g(x,u,y)=p(x,u)\,q(x,u,y)\,\chi(x,u,y)
          +\sigma(x,u,y)-C(x,u,y),                 \tag{1}
\]
where $\chi\in[0,1]$ is tariff collection, $\sigma$ is the declared subsidy,
and $C$ includes all operating, maintenance and replacement cash payments.
Investment resources are purchased at those cash costs; any transfer
component must be separated for social welfare. Donor exit is represented
by states in which $\sigma=0$. Known finite subsidies can be encoded by a
time state before exit. Cash $b$ becomes $b+g$ at the end of the period
and must never be negative. No expenditure is due earlier within a period;
if it is, divide that period into payment subdates.

A planned grid-arrival closure is permissible only if the specified
service obligation allows it. Such an action pays decommissioning and
terminal liabilities, enters an absorbing state with $g=0$, and cannot
silently discharge a continuing reliability obligation.

\section{An exact permanent-service reserve recursion}
Define the reserves required for zero further periods by $r_0(x)=0$ and
recursively
\[
 r_{n+1}(x)=\min_{u\in\mathcal U(x)}
       \max\left\{0,\max_{y\in\mathcal N(x,u)}
                         [r_n(y)-g(x,u,y)]\right\}.          \tag{2}
\]
The criterion is pathwise solvency for every permitted sequence of shocks,
not positive expected value or a high survival probability. This is useful
for a guarantee contract; less conservative probabilistic criteria are
different research questions.

\begin{theorem}[Minimum reserves for indefinite reliable operation]
For each $n$, $r_n(x)$ is the smallest initial cash allowing a
nonanticipative policy to meet all service and cash constraints for $n$
periods. The sequence is nondecreasing. Put
$r_\infty(x)=\sup_n r_n(x)\in[0,\infty]$.
A finite reserve $b$ permits indefinite operation from state $x$ if and
only if $b\geq r_\infty(x)$. Whenever $r_\infty(x)<\infty$, an action
can be selected depending only on the state such that all its successor
states have finite requirements and
\[
 r_\infty(y)\leq r_\infty(x)+g(x,u,y)
 \quad\text{for every permitted successor }y.              \tag{3}
\]
Selecting such actions at every finite-requirement state yields an
indefinitely feasible stationary policy.
\end{theorem}
\begin{proof}
For one additional period and fixed action $u$, feasibility requires
$b\geq0$ and $b+g(x,u,y)\geq r_n(y)$ for every successor, which is
exactly the inner maximum in (2). Finite action sets ensure that its
minimum is attained. Backward induction proves the finite-horizon
claim. The recursion is order-preserving and $r_1\geq r_0$, so induction
gives monotonicity. Every infinite feasible policy is feasible for every
finite horizon, proving necessity of $b\geq r_\infty(x)$.

Suppose $r_\infty(x)<\infty$. Choose a minimizing action at each horizon.
Some action $u$ occurs along an infinite subsequence, because the action
set is finite. On this subsequence (2) implies
$r_n(y)-g(x,u,y)\leq r_{n+1}(x)\leq r_\infty(x)$ for each successor.
Taking its unbounded-horizon limit yields (3), so every successor has a
finite requirement. Choose one such action for each finite-requirement
state. Starting with $b\geq r_\infty(x)$, (3) implies that next cash is
at least the next state's requirement and hence nonnegative. Induction
establishes this invariant for every finite date and every permitted shock
path, which is the infinite-horizon requirement. This proves sufficiency
without exchanging an infinite minimum and limit.
\end{proof}

Equation (2) is a verifiable engineering and finance certificate. It can
exclude a low tariff even when that tariff generates high usage, because
replacement reserves cannot accumulate. It can favor a maintenance
program if the added near-term cost reduces future replacement or outage
obligations. Actions here already impose reliability; the theorem does not
prove a physical reliability model from accounting alone.

\section{Closed-form renewal accounting}
Consider a deterministic plant that starts with a new battery purchased
as part of initial real cost $K_0$. At the end of each period it earns
collected operating margin $g=pq\chi-c$, with fixed reliable demand $q>0$,
collection $\chi>0$ and recurring resource cost $c$. At dates
$L,2L,\ldots$ it must pay $R>0$ to replace the battery. Replacement cost
is not included in $c$ or $K_0$ beyond the initial battery. Demand, tariff
and reliability are constant and satisfy the declared access obligations.

\begin{proposition}[Cash feasibility and full capital recovery]
With no subsequent subsidy and no borrowing, permanent operation is
possible with some finite cash reserve if and only if $Lg\geq R$.
When that condition holds the minimum reserve is zero under the stated
end-of-period payment convention. With discount $\delta\in(0,1)$,
full discounted project surplus is
\[
 \Pi=-K_0+\frac{g\delta}{1-\delta}
                 -\frac{R\delta^L}{1-\delta^L}.             \tag{4}
\]
Therefore subsidy-free operation with nonnegative project surplus and
an affordability cap $p\leq\bar p$ is possible precisely when a tariff
in the admissible interval satisfies
\[
 pq\chi-c\ \geq\
 \max\left\{\frac RL,
 \frac{1-\delta}{\delta}
       \left(K_0+\frac{R\delta^L}{1-\delta^L}\right)\right\}.\tag{5}
\]
Here $q,\chi,c$ are held fixed across the tariff interval; endogenous demand
requires solving (5) using its actual tariff-dependent margin instead.
\end{proposition}
\begin{proof}
At the end of date $t=kL+j$, $0\leq j<L$, cumulative cash from operations
and replacements is $k(Lg-R)+jg$. If $Lg<R$, it tends to negative infinity
along $j=0$, so every finite initial reserve fails eventually. If
$Lg\geq R$, then $g>0$ and every prefix is nonnegative, giving feasibility
from zero reserve. Discounted receipts and replacements are geometric
series, giving (4). Rearranging $\Pi\geq0$ gives the second lower bound
in (5), while permanent operation gives the first. Both conditions are
necessary and sufficient within the constant-margin model.
\end{proof}

For an arbitrary deterministic periodic cash sequence
$g_1,\ldots,g_L$ including replacements, the same argument gives a useful
extension: finite-reserve feasibility holds exactly when
$\sum_{j=1}^L g_j\geq0$, and its minimum reserve is
\[
 b^*=\max\left\{0,\max_{1\leq k\leq L}
                          -\sum_{j=1}^k g_j\right\}.         \tag{6}
\]
Indeed a negative cycle total eventually defeats every reserve; a
nonnegative cycle total makes the worst prefix occur in the first cycle.
Thus upfront maintenance can require reserves even when (5)'s simple
end-of-period convention does not.

\section{Why positive discounted value can coexist with certain failure}
Take $L=2$, $g=1$, $R=3$, $K_0=1/100$ and $\delta=1/10$.
Equation (4) gives
\[
 \Pi=-1/100+1/9-1/33>0,
\]
but each renewal cycle loses one unit of cash. Every finite reserve runs
out. This exact counterexample relies on no statistical estimation.
Discounting future losses heavily can create positive NPV while a
no-borrowing permanent-service requirement is impossible.

A stochastic version is equally restrictive for guarantees. Suppose one
state has a permitted self-loop producing $g=-d<0$ under every available
action. Repeating that self-loop for $n$ dates requires at least $nd$
cash, hence (2) implies $r_\infty=\infty$ there. If this adverse loop has
positive conditional probability each time, any finite reserve faces a
positive probability of failure, although expected period margin may be
positive. This does not say eventual failure has probability one; that
stronger assertion would require additional drift and recurrence facts.

\section{Economic use and unresolved questions}
For a specified finite action menu, the reserve theorem completely solves
the robust financing subproblem. One may compare governance regimes by
changing measured collection and transition primitives, not by assigning
ownership a favorable coefficient. An initial donor grant can cover $K_0$
and reserve $b^*$ but cannot cure a permanently negative cycle margin.
Ongoing subsidy changes (1) and must be recorded as a public cash outlay;
it is a transfer in a combined equal-weight welfare account, while battery,
fuel, maintenance and administration are real costs. Household benefits
are gross utility less these real costs, with tariff payments and operator
receipts canceled once. This distinguishes commercial return, supported
viability and social desirability.

The full target still requires credible tariff-dependent demand,
collection behavior, maintenance effectiveness, affordability thresholds,
transition probabilities and grid-arrival scenarios. A next empirical
step is to estimate a feasible action menu with confidence bounds on those
primitives, then report both reserve guarantees and a separately specified
probability-of-failure frontier. Finite-horizon operating profit without
those obligations cannot certify permanent service.

\begin{thebibliography}{9}
\bibitem{MM} R. Meeks and M. Mahadevan (2025), \emph{Electricity
Infrastructure}, VoxDevLit 15(1), May. Primary publisher abstract and
policy summary inspected, particularly cost recovery, reliability and
minigrid viability. \url{https://voxdev.org/voxdevlit/electricity-infrastructure}.
\end{thebibliography}

\end{document}
